\documentclass[11pt]{article}
\usepackage{amsmath,amssymb,amsthm}
\newtheorem{proposition}{Proposition}
\begin{document}
\title{Permutation covariance and the support problem}
\author{Research checkpoint}
\maketitle

This note studies a possible joint decoder, not the original worst-case
head complexity. It proves covariance identities and isolates gaps in a
probabilistic permutation argument. It supplies no improved upper bound.

\section{Exact covariance of a permutation orbit}

Let $H\geq2$. For each input $x$, let $A(x)$ be an $H$ by $H$ real matrix.
The entry $A_{hl}(x)$ is the contribution obtained by assigning level $l$
to fixed head direction $h$. For a uniform permutation $\pi$, put
\[
 S_\pi(x)=\sum_{h=1}^H A_{h,\pi(h)}(x).
\]
Write $\widetilde A(x)$ for the matrix after subtracting its row means and
column means and adding its grand mean.

\begin{proposition}
The exact covariance is
\[
 \operatorname{Cov}(S_\pi(x),S_\pi(y))
 =\frac1{H-1}\sum_{h,l=1}^H
       \widetilde A_{hl}(x)\widetilde A_{hl}(y).
\]
\end{proposition}
\begin{proof}
Double centering changes every permutation sum by the same constant.
For the centered matrices, the equal-row terms in the product expectation
contribute $H^{-1}\sum_{h,l}\widetilde A_{hl}(x)\widetilde A_{hl}(y)$.
The unequal-row terms contribute
\[
 \frac1{H(H-1)}\sum_{h,l}\widetilde A_{hl}(x)
       \sum_{g\ne h,\ r\ne l}\widetilde A_{gr}(y).
\]
The inner sum equals $\widetilde A_{hl}(y)$ because all row and column sums
vanish. Adding the two terms gives the formula.
\end{proof}

Thus an actual covariance computation requires double centering. Counting
$H!$ permutations or observing that their convex hull contains zero does
not prove that their evaluation vectors meet every strict sign orthant.
The latter distinction has an exact finite counterexample on the sibling
branch \texttt{ar/47c1ff45613b}, in the shared permutation dictionary audit.

\section{A reciprocal kernel on a Hamming layer}

Fix $m$-element subsets $x,y\subseteq[k]$, independent random signs
$v_i\in\{-1,1\}$, and $B>m\varepsilon>0$. Define
\[
 r_v(x)=\frac1{B+\varepsilon\sum_{i\in x}v_i}.
\]
For a fixed $j$-element subset $S\subseteq x$, let
\[
 c_j=\mathbb E_v\left[r_v(x)\prod_{i\in S}v_i\right].
\]
Permutation symmetry makes this depend only on $j$, not on $x$ or $S$.

\begin{proposition}
The exact second moment is
\[
 K(x,y)=\mathbb E_v[r_v(x)r_v(y)]
       =\sum_{j=0}^{|x\cap y|}\binom{|x\cap y|}{j}c_j^2.
\]
Moreover,
\[
 c_j=2^{-m}\sum_{a=0}^j\sum_{b=0}^{m-j}
    \frac{(-1)^a\binom ja\binom{m-j}b}
         {B+\varepsilon(m-2a-2b)}.
\]
\end{proposition}
\begin{proof}
Expand $r_v(x)$ in the orthonormal characters
$\prod_{i\in S}v_i$, indexed by $S\subseteq x$. Only identical characters
survive the expectation of the product. Their common supports lie in
$x\cap y$, giving the first formula. Counting negative signs in $S$ and
$x\setminus S$ gives the second formula.
\end{proof}

The Johnson harmonic subspace of degree $i$ has multiplicity
$\binom ki-\binom k{i-1}$. The standard inclusion-matrix decomposition
of the kernel $\binom{|x\cap y|}{j}$ gives its eigenvalue
\[
 \lambda_i=\sum_{j=i}^m c_j^2
       \binom{m-i}{j-i}\binom{k-j-i}{m-j}.
\]
This last spectral formula is recorded for calculation; the elementary
second-moment identity above does not require it.

As $B\downarrow m\varepsilon$, every $c_j$ has leading term
$(-1)^j/(2^m(B-m\varepsilon))$. Consequently the leading kernel is a
multiple of $2^{|x\cap y|}$, which is not the identity kernel. Approaching
the positivity boundary does not make the independent-sign ensemble
isotropic on a Hamming layer.

\section{Isotropy without a decoder}

There is a related ensemble whose limiting covariance is isotropic, but
its fixed samples reveal a support obstruction. Set $k=2m$ and choose an
$m$-element set $S$ uniformly. Define an admissible positive denominator
on the whole Boolean cube by
\[
 D_{S,\rho}(x)=\rho+\rho^2\sum_{i\in S}x_i
                         +\sum_{i\notin S}x_i,
 \qquad \rho>0.
\]
Every denominator coefficient is positive. On the $m$-th layer,
\[
 \lim_{\rho\downarrow0}\frac{\rho}{D_{S,\rho}(x)}
 =\mathbf1_{x=S}.
\]
Therefore the limiting second moment of a random head is exactly
$\binom{2m}m^{-1}$ times the identity. Nevertheless a fixed family of
$H$ such heads has at most $H$ supported layer points in this limit.
Permuting their scalar residues only changes amplitudes at those points.
It does not move support or give arbitrary signs on the remaining points.

The finite-$\rho$ functions have nonzero leakage elsewhere. Ignoring that
leakage is not a proof of a finite-$\rho$ impossibility theorem: it could
be amplified after cancellation. Conversely, using only the isotropic
average is not a proof that the leakage can encode independent labels.

\section{First unresolved estimate}

A successful random-orbit proof must establish strict orthant coverage
for the actual structured matrices $A_{hl}(x)$, after all allowed common
corrections have been charged. Replacing them by independent random
matrix entries is not justified. The Cauchy matrices arise from only
$O(Hk)$ fixed direction and numerator parameters, whereas an unrestricted
random array on $2^k$ inputs has $H^2 2^k$ independent entries.

Likewise, quotienting by a low-degree subspace requires an actual common
head realization of that quotient correction, with its number of heads
included. The covariance identity makes this quotient a definite linear
algebra problem but does not solve the realization or coverage problem.
No fixed-power or unbounded-factor improvement follows from this note.
\end{document}
